\documentclass[12pt,a4paper,english,twoside]{article}
\usepackage{fontspec}
\usepackage[english]{babel}
\usepackage{amsmath}
\usepackage{amsfonts}
\usepackage{amssymb}
\usepackage{graphicx}
\usepackage{fancyhdr}
\usepackage{lmodern}
\usepackage{minted}
\usepackage{hyperref}
\usepackage{verbatim}
\usepackage{alltt}
\usepackage{color}
\usepackage{listings}
\lstset{frame=single,keepspaces=true,numbers=none,tabsize=2}
\defaultfontfeatures{Ligatures=TeX}
\setmainfont{DejaVu Serif}
\setsansfont{DejaVu Sans}
\setmonofont{DejaVu Sans Mono}
\usepackage[cm]{fullpage}
\setlength{\headsep}{10pt}
\setlength{\footskip}{20pt}
\addtolength{\textheight}{-30pt}
\author{
    Bart de Boer \texttt{bart.de.boer@hightechict.nl}
    \and
    Bas Bossink \texttt{bas.bossink@hightechict.nl}
}
\date{\today}
\title{Functional Programming and Test Driven Development\\
\vspace{2 mm} {\large A match made in heaven}}
\newminted{scheme}{linenos=true,mathescape,xleftmargin=20px}
\usepackage{natbib}  \bibliographystyle{chicago}
\renewcommand{\headrulewidth}{0.5pt}
\renewcommand{\headheight}{35pt}
\fancyhead[L]{\tiny Functional Programming and Test Driven Development}
\fancyhead[R]{\includegraphics[scale=0.2]{HTLogo}}
\fancyfoot[C]{\thepage}
\pagestyle{fancy}
\begin{document}
\maketitle
\section{Scheme the language}
This section will briefly describe a small portion of the Scheme language that 
will be used in this workshop.
Before we start, a few typographical conventions, all text written in a 
\texttt{monospaced font} is either a Scheme keyword or text provided or 
returned from the Scheme interpreters REPL (Read Eval Print Loop), a console 
application that evaluates Scheme expressions as you provide them. The prompt 
of the REPL starts with a hash \texttt{\#}, followed by a semicolon \texttt{;}, the 
character that starts a comment in Scheme, followed by a count of the number 
of successfully evaluated expressions, followed by the greater than \texttt{>} 
symbol. For example \texttt{\#;37>}.
\subsection{Introduction}
Since Scheme is a LISP (LISt Processing) dialect the central syntactical 
element in the language is the list. But before we dive into lists lets look 
at Scheme's `primitive' expressions, that is expression that evaluate to 
themselves. In Scheme these are: booleans, numbers, characters, strings, 
symbols and vectors.
\subsection{Boolean values}
Scheme has the boolean values \texttt{\#f} and \texttt{\#t} for false and true 
respectively. In boolean expressions all values except \texttt{\#f} are 
treated as true.
\subsection{Numbers}
Scheme supports different types of numbers for this workshop we only need to 
use integers and floating point numbers in base 10.
\begin{itemize}
\item integers 
\begin{lstlisting}
#;7> 37
37
\end{lstlisting}
\item floating point values
\begin{lstlisting}
#;11> 1.3552e-12
1.3552e-12 
#;12> .2342
0.2342
\end{lstlisting}
\end{itemize}
\subsection{Strings}
Strings are written as a number of characters surrounded by double-quotes.
\begin{lstlisting}
#;18> "Fred Flinstone"
"Fred Flinstone"    
#;19> "Barney Rubbles"
"Barney Rubbles"
\end{lstlisting}
\subsection{Symbols}
A symbol is an immutable string for which strings with the same content 
will refer to the same object. Symbols can be compared with 
\texttt{eq?} (essentially pointer comparison), whereas strings need to be 
compared by looking at their contents with \texttt{equal?}. Let's look at some 
examples:
\begin{lstlisting}
#;20> 'Fred
Fred
#;22> (eq? 'Fred 'Fred)
#t
#;23> (eq? "Fred" "Fred")
#f
#;24> (equal? "Fred" "Fred")
#t
\end{lstlisting}
\subsection{Other types of objects}
\subsubsection{Lists}
A list is denoted as a \texttt{'('} followed by elements separated by spaces 
followed by a \texttt{')'}. Since the list is also the primary syntactical 
construct of the Scheme language lists are interpreted in a special way. The 
interpreter tries to evaluate the first element of the list to a procedure and 
applies that procedure with the rest of the elements in the list as the 
arguments to the procedure.
\begin{lstlisting}
#;22> (1 2 3)
Error: call of non-procedure: 1

Call history:

<syntax>	  (1 2 3)
<eval>	  (1 2 3)	<--

\end{lstlisting}
In the example output you can see the way Scheme interprets lists. The first 
element of the list is interpreted as a procedure if it is not one of the 
standard syntactic forms. This interpretation can be circumvented by 
\texttt{quote}-ing the expression using the \texttt{quote} syntactic form or 
it's abbreviated form \texttt{'}. An example of this is given below:
\begin{lstlisting}
#;22> (quote (1 2 3))
(1 2 3)
#;23> '(1 2 3)
(1 2 3)
#;24> (+ 1 2 3)
6
\end{lstlisting}
\paragraph{Creating lists}
Lists can be created in several different ways:
\begin{itemize}
\item List literals \begin{lstlisting}
#;25> (quote (1 2 3))
(1 2 3)
#;26> '(1 2 3)
(1 2 3)
\end{lstlisting}
\item Using the \texttt{list} procedure
\begin{lstlisting}
#;27> (list 1 2 3)
(1 2 3)
\end{lstlisting}
\item Using the \texttt{cons} procedure to construct lists  \begin{lstlisting}
#;28> (cons 1 (cons 2 (cons 3 '())))
(1 2 3)
\end{lstlisting}
The name \texttt{cons} stands for "construct". The \texttt{cons} procedure 
creates what is known as a cons cell, a pair of pointers. These pointer pairs 
can be used to create a single linked list where the second part of the pair 
points to the next element (= cons cell) in a list. The last cons cell of a 
list points to a distinguished value that is not a pair in our case the empty 
list \texttt{()}.  \end{itemize}
\paragraph{Deconstructing lists}
Lists can be deconstructed using the \texttt{car} and \texttt{cdr} procedures. 
These names stem from the original Lisp implementation in the late 1950's.  
\begin{itemize}
\item \texttt{car} short for "Contents of the Address part of Register number" 
  is referred to as \texttt{head} or \texttt{first} in some languages.
\item \texttt{cdr} short for "Contents of the Decrement part of Register 
  number" is referred to as \texttt{tail} or \texttt{rest} in some languages.
\end{itemize} \begin{lstlisting}
#;29> (car '(1 2 3))
1
#;30> (cdr '(1 2 3))
(2 3)
\end{lstlisting}
\subsubsection{Functions}
In Scheme procedures or functions are first class values as well. We'll look at 
them in detail later but first we are going to look at definitions.
\subsection{Bindings a.k.a Definitions}
To create a top-level binding the \texttt{define} syntactic form is used which 
in it's most basic form looks like this:
\begin{equation*}
\texttt{(define <variable> <expression>)}
\end{equation*}
Some examples:
\begin{lstlisting}
#;31> (define fred 37)
#;32> fred
37
#;33> (define barney "Betty")
#;34> barney
"Betty"
\end{lstlisting}
Note that \texttt{define} expressions can also appear inside an inner scope. 
We will see this later on during a couple of the exercises.
\subsection{Procedures/Functions}
\subsubsection{Creating procedures}
Procedures can be created using the \texttt{lambda} syntactical form. Here 
lambda refers back to the lambda calculus \cite{lambda}.
The basic syntax of a lambda expression looks like this:
\begin{equation*}
\texttt{(lambda (par$_{0} \dots$ par$_{n}$) <body>)}
\end{equation*}
Examples of \texttt{lambda} expressions:
\begin{lstlisting}
#;38> (lambda (x) (* x x))
#<procedure (? x)>
#;39> ((lambda (x) (* x x)) 2)
4 
#;40> ((lambda (x y) (* x y)) 3 4)
12
#;41> (define square (lambda (x) (* x x)))
#;43> (square 3)
9
\end{lstlisting}
From the last example you can take away that combining \texttt{define} with a 
\texttt{lambda} expression gives you a way to define functions. 
\subsubsection{Applying procedures}
We have already seen examples of applying a procedure to arguments above. 
There also is a handy function called \texttt{apply} that can be used to apply 
a procedure to it's arguments.
\begin{lstlisting}
#;44> (apply (lambda (x y) (* x y)) '(3 4))
12
#;45> ((lambda (x y) (* x y)) (car '(3 4)) (car (cdr '(3 4))))
12
\end{lstlisting} Using \texttt{apply} is sometimes handy because then you 
don't have to deconstruct the list as you can see in the example above.
\section{Fizzbuzz}
\subsection{Writing a solution for the basic Fizzbuzz kata}
Now that we have covered the most basic forms of Scheme we will look at some 
more Scheme syntactic forms using the Fizzbuzz kata as a vehicle. The result 
of completing this kata will be used in your next assignment.
\subsubsection{The assignment}
Write a function that for each positive integer $x$ will produce the following 
result:
\begin{enumerate}
  \item if $x$ is divisible by 3 return the string "fizz"
  \item if $x$ is divisible by 5 return the string "buzz"
  \item if $x$ is divisible by both 3 and 5 return the string "fizzbuzz"
  \item otherwise return $x$
\end{enumerate}
\subsection{Our first test}
Since this workshop is about Test Driven Development we will start with our 
first test. We will be using a Scheme `standard library' called SRFI 78. It is 
already provided for us in the check.scm file. To save some space in this 
document we will use the \texttt{report-failed} mode of our test library in 
the next runs. A first test would look like this:
\subsubsection{A first test case}
\begin{schemecode}
(load "~/.bin/check.scm")
(check-set-mode! 'report-failed)

(check (fizzbuzz 1) => 1)

(check-report)
\end{schemecode}
Our environment is configured such that our tests our automatically run when 
we save our Scheme source files. When this first test is run the following 
output is produced:
\begin{lstlisting}

Error: unbound variable: fizzbuzz

    Call history:

    <syntax>	  (##core#lambda () (fizzbuzz 1))
    <syntax>	  (##core#begin (fizzbuzz 1))
    <syntax>	  (fizzbuzz 1)
    <syntax>	  (##core#undefined)
    <eval>	  (>=426 check:mode427 1)
    <eval>	  (check:proc428 (quote429 (fizzbuzz 1) )...
    <eval>	  [check:proc] (eqv?60 tmp57 (quote 0))
    <eval>	  [check:proc] (eqv?60 tmp57 (quote 1))
    <eval>	  [check:proc] (eqv?60 tmp57 (quote 10))
    <eval>	  [check:proc] (eqv?60 tmp57 (quote 100))
    <eval>	  [check:proc] (check:report-expression expression)
    <eval>	  [check:report-expression] (newline)
    <eval>	  [check:report-expression] (check:write expression)
    <eval>	  [check:report-expression] (display " => ")
    <eval>	  [check:proc] (thunk)
    <eval>	  (fizzbuzz 1)	<--
\end{lstlisting}
Our test can not be interpreted because we have not yet defined our 
\texttt{fizzbuzz} procedure. Lets look at the code and explain each line 
one by one:

\begin{enumerate}
  \item line 1: loads our unit testing library of choice, the check library
  \item line 2: configures the test library to only report failures, for 
  brevitiy
  \item line 4: this is our first assert you could read this as: verify that calling 
    the \texttt{fizzbuzz} function with \texttt{1} will result in the value 
    \texttt{1}
 \item line 6: this last line creates a report of our results.
\end{enumerate}
\subsubsection{The first fizzbuzz implementation}
Now we have to define the first revision of our \texttt{fizzbuzz} function, 
to keep things simple we will write our production code in the same file as 
our test code. Here is the first implementation of our \texttt{fizzbuzz} 
function:
\begin{schemecode}
(load "~/.bin/check.scm")
(check-set-mode! 'report-failed)

(define fizzbuzz (lambda (x) 1))

(check (fizzbuzz 1) => 1)

(check-report)
\end{schemecode}
\begin{lstlisting}
; *** checks *** : 1 correct, 0 failed.
\end{lstlisting}
\subsubsection{The second test case}
Great our first passing test. Lets add our next test case.  
\begin{schemecode}
(load "~/.bin/check.scm")
(check-set-mode! 'report-failed)

(define fizzbuzz (lambda (x) 1))

(check (fizzbuzz 1) => 1)
(check (fizzbuzz 2) => 2)

(check-report)
\end{schemecode}
\begin{lstlisting}
(fizzbuzz 2) => 1 ; *** failed ***
 ; expected result: 2

; *** checks *** : 1 correct, 1 failed. First failed example:

(fizzbuzz 2) => 1 ; *** failed ***
 ; expected result: 2
\end{lstlisting}
\subsubsection{Making our second test case pass}
Making this failing test pass is a 
matter of replacing our constant with a parameter.
\begin{schemecode}
(load "~/.bin/check.scm")
(check-set-mode! 'report-failed)

(define fizzbuzz (lambda (x) x))

(check (fizzbuzz 1) => 1)
(check (fizzbuzz 2) => 2)

(check-report)
\end{schemecode}
\begin{lstlisting}
; *** checks *** : 2 correct, 0 failed.  
\end{lstlisting}
Excellent. Now lets refactor our (test) code a bit to make adding test cases 
a bit simpler.
\section{Assignment 1: Remove duplication from our test cases}
\subsection{The work environment}
When you start the provided virtual machine you are logged in automatically 
and are presented with a split (tmux) screen. The top pane asks you which 
editor you would like to use during this session. If you are not familiar with 
any of the listed editors, we recommend you use Diakonos, it is an easy to use 
relatively feature rich console editor. A list with keyboard shortcuts is 
provided for easy reference. The bottom pane is empty, it is running a 
\texttt{watchr} script that will automatically run the tests when they are 
saved.  
\subsection{Writing the \texttt{check-fizzbuzz} function}
The code for this exercise is in the \texttt{fizzbuzz.scm} file. By choosing 
any of the available editors from the menu, it will be opened for you 
automatically. Write a function \texttt{check-fizzbuzz} that takes as 
parameters the value for which we want to check fizzbuzz and the expected 
outcome. The function should call the check function with the result of 
applying fizzbuzz and the expected value. Replace the two test cases with 
calls to the function you just created. You will get 5 minutes for this small 
exercise after which we'll discuss the solution and continue with implementing 
the kata. So don't turn the page just yet. Ahh, you've finished the assignment 
turn the page and check your solution.
\vfill
\pagebreak
\subsection{Solution for assignment 1}
\begin{schemecode}
(load "~/.bin/check.scm")
(check-set-mode! 'report-failed)

(define fizzbuzz (lambda (x) x))

(define check-fizzbuzz 
  (lambda (value-to-check expected)
    (check (fizzbuzz value-to-check) => expected)))

(check-fizzbuzz 1 1)
(check-fizzbuzz 2 2)

(check-report)
\end{schemecode}
\begin{lstlisting}
; *** checks *** : 2 correct, 0 failed.  
\end{lstlisting}

\subsubsection{The third test case}
\begin{schemecode}
(load "~/.bin/check.scm")
(check-set-mode! 'report-failed)

(define fizzbuzz (lambda (x) x))

(define check-fizzbuzz 
  (lambda (value-to-check expected)
    (check (fizzbuzz value-to-check) => expected)))

(check-fizzbuzz 1 1)
(check-fizzbuzz 2 2)
(check-fizzbuzz 3 "fizz")

(check-report)
\end{schemecode}
\begin{lstlisting}
(fizzbuzz 3) => 3 ; *** failed ***
 ; expected result: "fizz"

; *** checks *** : 2 correct, 1 failed. First failed example:

(fizzbuzz 3) => 3 ; *** failed ***
 ; expected result: "fizz"
\end{lstlisting}

To be able to make the failing test pass we need to know how we can write 
conditional expressions in Scheme, luckily that is the topic of the next 
section.

\subsection{Conditionals}
\paragraph{\texttt{if}}
The simplest conditional expression in Scheme is the \texttt{if} expression 
which
has the following syntax:
\begin{equation*}
    \texttt{(if <expression> <then expression> [<else expression>])}
\end{equation*}
The value of the if expression is the value of the expression that gets 
evaluated depending on the value of the \texttt{<expression>}, the else 
expression is optional. If the \texttt{<expression>} evaluates to false and 
the else expression is missing than the value of the if expression is 
undefined.

Now we can make our failing test pass:
\subsubsection{Making the third test case pass}
\begin{schemecode}
(load "~/.bin/check.scm")
(check-set-mode! 'report-failed)

(define fizzbuzz 
  (lambda (x) 
    (if (= 3 x) "fizz" x)))

(define check-fizzbuzz 
  (lambda (value-to-check expected)
    (check (fizzbuzz value-to-check) => expected)))

(check-fizzbuzz 1 1)
(check-fizzbuzz 2 2)
(check-fizzbuzz 3 "fizz")

(check-report)
\end{schemecode}
\begin{lstlisting}
; *** checks *** : 3 correct, 0 failed.
\end{lstlisting}
The next interesting case is of course the number 5.  
\subsubsection{The fourth test case}
\begin{schemecode}
(load "~/.bin/check.scm")
(check-set-mode! 'report-failed)

(define fizzbuzz 
  (lambda (x) 
    (if (= 3 x) "fizz" x)))

(define check-fizzbuzz 
  (lambda (value-to-check expected)
    (check (fizzbuzz value-to-check) => expected)))

(check-fizzbuzz 1 1)
(check-fizzbuzz 2 2)
(check-fizzbuzz 3 "fizz")
(check-fizzbuzz 5 "buzz")

(check-report)
\end{schemecode}
To make it pass we could use a nested if but Scheme provides us with the 
\texttt{cond} expression which reads a lot nicer.
\paragraph{\texttt{cond}}
Below the syntax for the \texttt{cond} syntactic form:
\begin{lstlisting}[mathescape]
  (cond 
    ((<predicate>$_{0}$) <expression>$_{0}$)
    $\vdots$ 
    ((<predicate$_{n}$) <expression>$_{n}$)
    [(else <expression$_{n+1}$)])
\end{lstlisting}
Each of the predicates gets evaluated in order of appearance, the value of the 
\texttt{cond} expression is the value of the expression of the first predicate 
that evaluated to true. A \texttt{cond} expression can have an optional 
\texttt{else} clause that will be evaluated if none of the predicates 
evaluated to true.
\subsubsection{Making the fourth test case pass}
\begin{schemecode}
(load "~/.bin/check.scm")
(check-set-mode! 'report-failed)

(define fizzbuzz 
  (lambda (x) 
    (cond ((= 3 x) "fizz")
          ((= 5 x) "buzz")
          (else x))))

(define check-fizzbuzz 
  (lambda (value-to-check expected)
    (check (fizzbuzz value-to-check) => expected)))

(check-fizzbuzz 1 1)
(check-fizzbuzz 2 2)
(check-fizzbuzz 3 "fizz")
(check-fizzbuzz 5 "buzz")

(check-report)
\end{schemecode}
\begin{lstlisting}
; *** checks *** : 4 correct, 0 failed.  
\end{lstlisting}
\section{Assignment 2: Finish the basic fizzbuzz kata}
\subsection{The fifth test case}
The rest of the kata can now be implemented, all the Scheme constructs that we 
will need have been introduced. All that we need to know is, does Scheme have 
something like a `remainder' or `modulus' function. The function to use here 
is the \texttt{modulo} function. Let's add another test case that will guide 
us to using this \texttt{modulo} function. The code for this exercise is in 
the \texttt{fizzbuzz-second.scm} file. Open it using the editor you choose at 
the start of the session. Write the next test case and make it pass using the 
before mentioned \texttt{modulo} function. Optionally you can try out the 
behavior of the \texttt{modulo} function by switching to the next tmux window 
by pressing \texttt{Ctrl-a n}, and using the interpreter that is already 
started their for you. You can switch back to the editor window by pressing 
\texttt{Ctrl-a p}.
\subsection{Add the next test and make it pass}
Your next assignment is to add the next test case and make it pass. When all 
your tests pass you can move on to the next assignment.
\subsection{Introduce a helper function}
\subsubsection{Writing \texttt{divisible-by?}}
Up to now we have neglected to refactor the code 
we've written, when all tests were passing. Your next assignment is to define 
a \texttt{divisible-by?} function. That can be used like in the example below:
\begin{lstlisting}
;#46> (divisible-by? 3 6)
#t
\end{lstlisting}
\subsubsection{Call the new \texttt{divisible-by?} function from 
\texttt{fizzbuzz}}
Replace the predicates in the \texttt{cond} with applications of our new 
\texttt{divisible-by?} function and make sure all tests still pass.
\subsection{Finish the \texttt{fizzbuzz} implementation}
Finish the \texttt{fizzbuzz} implementation, start by adding a test case, 
watch it fail and then make it pass. 
\subsection{Reduce duplication further using \texttt{map}}
Now we will remove even more duplication by using the \texttt{map} and 
\texttt{apply} functions. We saw the \texttt{apply} function in the 
introduction it takes a function and a list as arguments and applies the 
supplied function to the supplied arguments in the list. The \texttt{map} 
function takes a function and a list as its parameters and returns a new list 
consisting of the function applied to each of the elements in the list.
Here is a short example
\begin{lstlisting}
#;46> (map (lambda (x) (* x x)) '(1 2 3))
(1 4 9)
\end{lstlisting}
The assignment is to define a list of test cases and run them using the 
\texttt{map} and \texttt{apply} functions.
\section{Assignment 3: Extending the rules of the fizzbuzz kata}
The rules for the fizzbuzz kata can be extended like so
\begin{enumerate}
  \item if $x$ is divisible by 3 return the string "fizz"
  \item if $x$ is divisible by 5 return the string "buzz"
  \item if $x$ is divisible by both 3 and 5 return the string "fizzbuzz"
  \item if the sequence of digits that denote $x$ contains 3 return \texttt{"fizz"} 
  \item if the sequence of digits that denote $x$ contains 5 return \texttt{"buzz"} 
  \item if the sequence of digits that denote $x$ contains both 3 and 5 return 
    \texttt{"fizzbuzz"} 
  \item otherwise return $x$
\end{enumerate}
Adapt the \texttt{fizzbuzz} function such that it adheres to this new set of 
rules. Do this by following the red, green, refactor cycle.
\begin{enumerate}
  \item Write a function \texttt{digits} that returns the list of digits (as 
    integers) that is used to denote the input integer. Recall from 
    mathematics that numbers in the decimal system are denoted as a sum of 
    powers of 10. E.g. $137 = 1 \times 10^{2} + 3 \times 10^{1} + 7 \times 
    10^{0}$.
  \item Write a function \texttt{contains?} that takes an element and a list 
    as it's parameters and returns \texttt{\#t} if the list contains the given 
    element and \texttt{\#f} otherwise.
  \item Combine the \texttt{digits} and \texttt{contains?} functions into a 
    \texttt{contains-digit?} function that takes a digit (an integer $x$ with 
    $0 \leq x < 10$) and an integer and returns \texttt{\#t} if the digit is 
    part of the decimal notation of the given number and \texttt{\#f} 
    otherwise.
  \item Finish the fizzbuzz kata with the extended rules using the helper 
    functions you've just created.
\end{enumerate}
\vfill
\pagebreak
\section{Assignment 4: Legacy code}
Some of you may know the legendary Hewlett-Packard calculators from the good 
old days that use Reverse Polish Notation (RPN). A form of notation that first 
denotes the operands followed by the operator. This assignment works with an
RPN calculator written in Scheme. It uses a stack to hold intermediary 
results, it can read a Scheme expression pushes all numbers it encounters onto 
the stack and when it encounters an operator pops the 2 operands of the stack 
applies the operator and pushes the result back onto the stack. Below you find 
some examples of using this calculator.
\begin{lstlisting}
#;47>(load "stack-calculator.scm")
#;48>(rpn '(37 42 23 + *))
2405
\end{lstlisting}
Lets first briefly look over the code:
\begin{schemecode}
(load "~/.bin/check.scm")
(check-set-mode! 'report-failed)

(define operator
  (lambda (expression) 
    (car expression)))

(define operator? 
  (lambda (expression)
    (define token (car expression))
    (or (eq? '+ token)
        (eq? '- token)
        (eq? '* token)
        (eq? '/ token))))

(define calculate 
  (lambda (expr stack)
    (cond ((null? expr) (car stack))
          ((operator? expr)
            (define first-operand (cadr stack))
            (define second-operand (car stack))
            (define result (eval (list 
                                   (operator expr) 
                                   first-operand 
                                   second-operand)))
            (calculate (cdr expr) (cons result (cddr stack))))
          (else (calculate (cdr expr) (cons (car expr) stack))))))

(define (rpn expression)
  (if (null? expression)
    "there is nothing to calculate!"
    (calculate expression '())))

(check (rpn '()) => )

(check-report)
\end{schemecode}
This code contains a few functions that we have not yet discussed:
\begin{enumerate}
\item line 11: \texttt{or} returns the first non \texttt{\#f} argument of the 
  supplied arguments without evaluating the rest of the arguments
  \item line 18: \texttt{null?} returns true if its parameter is the empty list
  \item line 20: \texttt{cadr} is a composition of \texttt{cdr} and \texttt{car} 
    \begin{equation*}
      \texttt{(cadr alist)} \Rightarrow \texttt{(car (cdr list))}
    \end{equation*}
  \item line 22: \texttt{eval} evaluates the given list in the same manner as the 
    Scheme interpreter would, which means evaluate the first element of the 
    list to a procedure definition evaluate the rest of the elements and 
    supply them as the formal parameters to the procedure.
  \item line 26: \texttt{cddr} is a composition of \texttt{cdr} with itself
    \begin{equation*}
      \texttt{(cddr alist)} \Rightarrow \texttt{(cdr (cdr alist))}
    \end{equation*}
\end{enumerate}
\subsection{Building the safety net}
Start by writing some tests that validate the current behavior of the 
\texttt{rpn} function. Open the \texttt{stack-calculotor.scm} file and add 
some tests to it.When you feel you have captured enough of the current 
behavior of the function proceed to the next assignment.
\subsection{Adding the \texttt{\^{}} operator}
Extend the \texttt{rpn} function such that it can understand expressions that 
contain the \texttt{\^{}} operator which should calculate the power function.  
\begin{equation*}
\texttt{(rpn '(2 3 \^{}))} \Rightarrow 2^{3} \Rightarrow 8 
\end{equation*}
The Scheme function for exponentiation is \texttt{expt}.
\begin{lstlisting}
#;48> (expt 2 3)
8
\end{lstlisting}
\begin{enumerate}
  \item Write a few test cases for the \texttt{operator?} function.
  \item Write a few test cases for the \texttt{operator} function.
  \item Add support for the \texttt{\^{}} operator to the \texttt{operator?} function.
  \item Add support for the \texttt{\^{}} operator to the \texttt{operator} function. 
  \item Add support for the \texttt{\^{}} operator to the \texttt{calculate} function.
\end{enumerate}
\subsection{Adding the \texttt{sqrt} function}
Up to now the calculator only understands binary operators, extend the 
calculator such that expressions can contain the \texttt{sqrt} operator that 
applies the \texttt{sqrt} function to the operand on the top of the stack.  
\begin{equation*}
  \texttt{(rpn '(9 sqrt))} \Rightarrow \sqrt{9} \Rightarrow 3
\end{equation*}
The Scheme function that calculates the square root of its argument is called 
\texttt{sqrt}.
\begin{enumerate}
  \item Create the \texttt{unary-operator?} function that can tell if the 
    first element in the expression is a unary operator or not. 
  \item Use this function in the \texttt{operator?} function.
  \item Extend the \texttt{operator} function to return the \texttt{sqrt} 
    function as a translation for \texttt{sqrt}.
  \item Extend the \texttt{calculate} function to take unary operators into 
    account.
  \item Introduce an \texttt{evaluate} function that calls \texttt{eval}. You 
    need to make this a function with a variable number of arguments. Such a 
    function can be defined using the following syntax
    \begin{equation*}
      \texttt{(lambda (<parameter>$_{0} \dots $<parameter>$_{n}$ . <rest>) <body>)}
    \end{equation*}
    In the body of the \texttt{lambda} expression the \texttt{rest} parameter 
    will be bound to a list containing the parameters$_{n+1}$ and beyond.
  \item Use this evaluate function to simplify the \texttt{calculate} 
    function.
\end{enumerate}
\section{Optional: add support for the Greatest Common Divisor}
The greatest common divisor (GCD) of two integers a,b is the largest integer that 
divides both a and b without leaving a remainder. For example: 
\begin{lstlisting}
#;49> (gcd 2 3)
1
#;50> (gcd 2 4)
2
#;51> (gcd 12 9)
3
#;52> (gcd 12 20)
4
\end{lstlisting}
Around 300 B.C Euclid described an algorithm for calculating the GCD that is 
still in use today. Wikipedia \cite{euclid} describes the algorithm  as follows:
\begin{quote}
 The Euclidean algorithm proceeds in a series of steps such that the output of 
 each step is used as an input for the next one. Let $k$ be an integer that 
 counts the steps of the algorithm, starting with zero. Thus, the initial step 
 corresponds to $k = 0$, the next step corresponds to $k = 1$, and so on.

 Each step begins with two nonnegative remainders $r_{k-1}$ and $r_{k-2}$. Since 
 the algorithm ensures that the remainders decrease steadily with every step, 
 $r_{k-1}$ is less than its predecessor $r_{k-2}$. The goal of the $k^{th}$ step is to find a 
 quotient $q_{k}$ and remainder $r_{k}$ such that the equation is satisfied
 \begin{equation*}
   r_{k-2} = q_{k} \times r_{k-1} + r_{k}
\end{equation*}
where $r_{k} < r_{k-1}$. In other words, multiples of the smaller number 
$r_{k-1}$ are subtracted from the larger number $r_{k-2}$ until the remainder 
is smaller than the $r_{k-1}$.

In the initial step $(k = 0)$, the remainders $r_{-2}$ and $r_{-1}$ equal $a$ and $b$, the 
numbers for which the GCD is sought. In the next step $(k = 1)$, the remainders 
equal $b$ and the remainder $r_{0}$ of the initial step, and so on. Thus, the 
algorithm can be written as a sequence of equations
\begin{eqnarray*}
  a &=& q_{0} \times b + r_{0} \\
  b &=& q_{1} \times r_{0} + r_{1} \\
  r_{0} &=& q_{2} \times  r_{1} + r_{2} \\
  r_{1} &=& q_{3} \times r_{2} + r_{3} \\
\vdots
\end{eqnarray*}
If $a$ is smaller than $b$, the first step of the algorithm swaps the numbers. For 
example, if $a < b$, the initial quotient $q_{0}$ equals zero, and the remainder $r_{0}$ 
is $a$. Thus, $r_{k}$ is smaller than its predecessor $r_{k-1}$ for all $k \geq 0$.

Since the remainders decrease with every step but can never be negative, a 
remainder $r_{N}$ must eventually equal zero, at which point the algorithm 
stops. The final nonzero remainder $r_{N-1}$ is the greatest common divisor of 
$a$ and $b$. The number $N$ cannot be infinite because there are only a finite 
number of nonnegative integers between the initial remainder $r_{0}$ and zero. 
\end{quote}
\subsection{Write the \texttt{gcd} function}
Write a (recursive) function \texttt{gcd} that calculates the GCD of it's two 
parameters $a$ and $b$. For simplicity you may assume that $b$ is smaller than 
$a$.
\subsection{Add support for the \texttt{gcd} function to the calculator}
Use the function created in the previous assignment to extend the calculator 
with a \texttt{gcd} function.
\bibliography{handout}
\end{document}
